\begin{figure}[htbp]
  \centering
  \begin{tikzpicture}[x=1cm, y=1cm, font=\scriptsize,
      eq/.style={ntcaixa, text width=1.7cm, minimum height=9mm},
      med/.style={circle, draw=corCinza, fill=white, line width=.5pt,
                  inner sep=0pt, minimum size=3.1mm, font=\tiny}]
    % ------------------------------------------------ fronteira -------
    \begin{scope}[on background layer]
      \fill[corFundo, rounded corners=3pt] (0,-1.35) rectangle (11.85,1.55);
      \draw[corTinta, line width=.7pt, dash pattern=on 3pt off 2pt,
            rounded corners=3pt] (0,-1.35) rectangle (11.85,1.55);
    \end{scope}
    \node[ntrot, anchor=north west, text=corTinta] at (0.15,1.45)
      {\textbf{Fronteira do indicador}};

    % ---------------------------------------------- cadeia processual --
    \node[eq] (pa1) at (1.75,0.30) {Pré-aquecimento};
    \node[eq] (des) at (3.90,0.30) {Dessalgação};
    \node[eq] (pa2) at (6.05,0.30) {Pré-aquecimento};
    \node[eq, draw=corFG, line width=.7pt, fill=corFG!7] (for) at (8.20,0.30)
      {Forno atmosférico};
    \node[ntcaixa, text width=1.6cm, minimum height=17mm,
          draw=corFG, line width=.7pt, fill=corFG!7] (col) at (10.70,0.20)
      {Coluna de destilação};

    \node[ntmini, anchor=north west] at (0.12,0.22) {crude};
    \draw[ntsetaf] (0.15,0.30) -- (pa1.west);
    \draw[ntsetaf] (pa1) -- (des);
    \draw[ntsetaf] (des) -- (pa2);
    \draw[ntsetaf] (pa2) -- (for);
    \draw[ntsetaf] (for) -- (col.west|-for);

    % ----------------------------------------------------- produtos ----
    \foreach \y/\t in {0.90/{gases e nafta}, 0.20/{petróleo e gasóleos},
                       -0.50/{resíduo atmosférico}} {
      \draw[ntsetaf] (11.50,\y) -- (12.15,\y);
      \node[ntmini, anchor=west, text width=2.5cm, align=left] at (12.20,\y) {\t};
    }

    % ------------------------------------------- entradas de energia ---
    % fuel gas: entra no forno, que e onde e queimado
    \draw[-{Stealth[length=4pt,width=3pt]}, draw=corFG, line width=.9pt]
      (8.20,-2.02) -- (for.south);
    \node[ntmini, anchor=north, text=corFG] at (8.20,-2.14) {fuel gás};
    \node[med] at (8.20,-1.35) {M};
    % eletricidade e vapores de alta: atravessam a fronteira, distribuidos
    % por varios equipamentos; a seta para na fronteira, sem atribuir destino
    \foreach \x/\c/\t in {1.75/corEE/{eletricidade}, 3.90/corV24/{vapor 24 e 10 bar}} {
      \draw[-{Stealth[length=4pt,width=3pt]}, draw=\c, line width=.9pt]
        (\x,-2.02) -- (\x,-1.20);
      \node[ntmini, anchor=north, text=\c] at (\x,-2.14) {\t};
      \node[med] at (\x,-1.35) {M};
    }
    \node[ntmini, anchor=north west, align=left] at (0.05,-2.50)
      {distribuídos por vários equipamentos; a eletricidade é medida ao mês e
       alocada ao dia};
    \draw[-{Stealth[length=4pt,width=3pt]}, draw=corV3, line width=.9pt]
      (10.70,-2.02) -- (col.south);
    \node[ntmini, anchor=north, text=corV3] at (10.70,-2.14) {vapor 3 bar};
    \node[med] at (10.70,-1.35) {M};

    % -------------------------------------- retornos e producao -------
    \draw[{Stealth[length=4pt,width=3pt]}-, draw=corV10, line width=.9pt]
      (6.05,2.00) -- (6.05,1.20);
    \node[ntmini, anchor=south, text=corV10] at (6.05,2.03)
      {vapor produzido na recuperação de calor};
    \node[med] at (6.05,1.55) {M};
    \draw[{Stealth[length=4pt,width=3pt]}-, draw=corCinza, line width=.7pt]
      (10.70,2.00) -- (col.north);
    \node[ntmini, anchor=south] at (10.70,2.03) {condensados};

    % ----------------------------------------------------- legenda -----
    \node[med] at (0.10,-3.20) {M};
    \node[ntmini, anchor=west] at (0.32,-3.20) {ponto de medição declarado};
    \draw[-{Stealth[length=4pt,width=3pt]}, draw=corCinza, line width=.9pt]
      (3.7,-3.20) -- ++(0.5,0);
    \node[ntmini, anchor=west] at (4.3,-3.20) {consumo bruto que entra};
    \draw[{Stealth[length=4pt,width=3pt]}-, draw=corV10, line width=.9pt]
      (7.6,-3.20) -- ++(0.5,0);
    \node[ntmini, anchor=west] at (8.2,-3.20) {produção creditada (abate ao bruto)};
  \end{tikzpicture}
  \caption[Fronteira energética da unidade de destilação atmosférica]%
  {Esquema funcional da unidade de destilação atmosférica e da fronteira
   que o indicador de desempenho usa. O desenho não tem escala nem caudais:
   representa a sequência de operações e os pontos onde entra e sai energia.
   Três noções de consumo convivem e não devem ser confundidas: o
   \emph{consumo bruto} é tudo o que atravessa a fronteira para dentro; o
   \emph{consumo líquido} desconta a produção creditada, como o vapor gerado
   na recuperação de calor; o \emph{consumo alocado} reparte por unidade
   aquilo que é medido num ponto comum, segundo uma regra declarada --- é
   o caso da eletricidade, medida ao mês para um conjunto de equipamentos e
   repartida por dia. Por isso, um desvio medido nesta fronteira pode ter
   origem no processo ou na regra de repartição.}
  \label{fig:fronteira-cdu}
\end{figure}
